\scanpage{41}
\[
 H^*(X)_T=H^*(X)\otimes_K\widehat K_T.
\]
une décomposition
\[
 \boxed{\varphi_T:\coprod_{p,q}H^q(X,\Omega_X^p)_T\longrightarrow H^*(X)_T}
\]
Cela permet de définir
\[
 H_T^{p,q}(X)=\varphi_T(H^q(X,\Omega_X^p)_T)\cap H^*(X)\subset F^p(H^q(X))
\]
qui est un \uline{sous-\inserted{$K$-}espace vectoriel de} $H^*(X)$, défini canoniquement par la donnée de $T$. Nous poserons de \unsure{$=$}
\[
 H_T^{pq}(X,\mathbf R)=\varphi_T(H^q(X,\Omega_X^p)_T)\cap H^*(X_{\unsure{\widehat K_T}},\mathbf R)\cap H^*(X)
\]
\deleted{c'est un \unsure{sous-} \deleted{$\unsure{K}$} $\unsure{\mathbf R,K}$-espace vectoriel de}
c'est encore un sous-\inserted{$K$-}espace vectoriel de $H^*(X)$ si $\widehat K_T\simeq\mathbf R$; \deleted{sinon} \inserted{$\widehat K_T\simeq\mathbf C$} c'est un sous-\deleted{$\unsure{R_{\mathrm{top}}}$}$R_T$-espace vectoriel, où
\[
 R_T=\mathbf R\cap K,\qquad\text{dans }\widehat K_T\simeq\mathbf C
\]
[sous-\uline{corps réel} pour $T$]; à fortiori, \deleted{\unsure{est}} dans tous les cas, on obtient un $K_r$-espace vectoriel, où
\[
 K_r\simeq\bigcap_T R_T\simeq\text{plus gd sous-corps tot. réel de }K
\]
(NB $K_r$ est une ext. algébrique du corps premier $\mathbf Q$).
\mspage{40}
